\section{Обзор существующих решений}
\label{sec:Chapter1} \index{Chapter1}

В данной главе рассматриваются два класса альтернатив: каналы передачи
диагностической информации из гостевой системы на хост и способы организации
буфера внутри такого канала. Эти вопросы связаны, но не совпадают: например,
\texttt{virtio-serial} является каналом передачи, а eBPF ring buffer ---
вариантом организации очереди событий.

\subsection{Последовательный порт и debugcon}
\label{sec:ch1-serial}

Наиболее простой способ вывести диагностическое сообщение из гостевой системы ---
использовать последовательный порт или отладочную консоль
(\texttt{debugcon})~\cite{qemu_debugcon}. Такой канал доступен на ранних этапах
загрузки и хорошо поддерживается прошивками: стандартная \texttt{DebugLib} в EDK2
как правило выводит сообщения именно в serial-порт.

Недостаток этого подхода связан с тем, что запись выполняется через операции
ввода-вывода или эмулируемые MMIO-регистры. В виртуальной машине такое обращение
обычно приводит к выходу из гостевого режима в гипервизор, обработке операции в
эмуляторе и возврату управления гостю. Если сообщение передаётся побайтно, то
эти накладные расходы возникают для каждого символа. Поэтому последовательный
порт удобен как универсальный и простой механизм, но плохо подходит для
интенсивного потока диагностических событий.

Существуют способы уменьшить накладные расходы. Например, доступ к портам
ввода-вывода можно разрешить пользовательскому процессу через TSS и карту
разрешения портов ввода-вывода, а сами сообщения можно складывать в
дополнительный кольцевой буфер на стороне гостя. Однако в таком случае
получается отдельный протокол поверх последовательного устройства. Для него
нужно отдельно решать задачи обнаружения канала в UEFI и Windows, передачи
метаинформации о буфере, синхронизации нескольких писателей и объединения
сообщений с логом эмулятора.

\subsection{virtio-console и virtio-serial}
\label{sec:ch1-virtio}

Устройства семейства virtio предназначены для эффективного взаимодействия гостя
и хоста в виртуализированной среде~\cite{virtio_spec}. Для передачи текстовых
данных могут использоваться \texttt{virtio-console} и
\texttt{virtio-serial}~\cite{virtio_console}. Они применяют virtqueue и лучше
масштабируются, чем побайтовая запись в последовательный порт.

Тем не менее для рассматриваемой задачи у virtio-устройств есть ограничения:
\begin{itemize}
    \item полноценная работа virtio требует инициализации устройства и драйвера,
    поэтому канал неудобен для самых ранних стадий UEFI;
    \item запись диагностического сообщения проходит через общий стек virtio, что
    добавляет служебные структуры, уведомления и обработку очередей;
    \item пользовательское приложение не может безопасно и напрямую писать в
    virtqueue без участия драйвера и согласованного протокола доступа;
    \item интеграция сообщений гостевой ОС и сообщений самого эмулятора требует
    отдельной логики синхронизации и форматирования.
\end{itemize}

Таким образом, virtio-console и virtio-serial являются альтернативой для уже
загруженной ОС, но хуже соответствуют требованиям раннего логирования,
потенциального отображения буфера в пользовательское пространство и минимизации
накладных расходов на запись.

\begin{table}[H]
    \centering
    \caption{Сравнение каналов передачи}
    \label{tab:channels-comparison}
    \small
    \begin{tabularx}{\textwidth}{|p{3.2cm}|X|X|}
        \hline
        Канал & Преимущества & Ограничения \\
        \hline
        Serial/debugcon &
        Простота, доступность на ранней загрузке &
        Частые VM-exit при побайтовой записи, слабая масштабируемость \\
        \hline
        virtio-console/virtio-serial &
        Стандартный канал виртуализации, virtqueue вместо побайтового вывода &
        Требует инициализации драйвера, не даёт прямой записи из userspace в общий буфер \\
        \hline
        LogDev &
        Разделяемая память, ранний синхронный режим, отдельное уведомление о появлении данных &
        Требует собственного виртуального устройства и драйверов \\
        \hline
    \end{tabularx}
\end{table}

\subsection{Организация кольцевых буферов}
\label{sec:ch1-ringbuf-alternatives}

После выбора канала на основе разделяемой памяти остаётся задача организации
самого буфера. Для передачи потока сообщений можно использовать несколько типов
очередей:
\begin{enumerate}
    \item \textbf{SPSC-буфер} (Single Producer Single Consumer) прост и быстр,
    если источник сообщений ровно один. Его можно использовать и с несколькими
    писателями, добавив внешний глобальный мьютекс, но тогда конкурентная запись
    фактически сериализуется, а стоимость вызова записи начинает зависеть от
    владельца блокировки.
    \item \textbf{MPMC-очередь} (Multiple Producer Multiple Consumer) поддерживает
    несколько писателей и читателей, но для данной задачи избыточна: читателем
    является только поток на стороне эмулятора.
    \item \textbf{MPSC-буфер с блокировкой} соответствует числу участников
    протокола, но блокировка делает задержку записи зависимой от планирования
    потоков внутри гостя и от планирования виртуальных CPU на хосте.
    \item \textbf{MPSC-буфер с атомарным резервированием} также соответствует
    модели «несколько писателей --- один читатель». Он требует более сложного
    алгоритма резервирования и фиксации записи, но уменьшает объём синхронизации
    в критическом пути записи.
\end{enumerate}

В Linux eBPF используется кольцевой буфер для передачи событий из программ eBPF
в пользовательское пространство~\cite{ebpf_ringbuf}. Он решает похожую задачу:
несколько источников событий формируют сообщения, а пользовательский процесс
считывает их из общей области памяти. Поэтому eBPF ring buffer рассматривается
в работе как алгоритмическая альтернатива для организации очереди событий. Его
ключевое отличие состоит в том, что этап резервирования места синхронизируется
с помощью блокировки.

В данной работе выбран MPSC-буфер с атомарным резервированием, так как он
соответствует модели взаимодействия «несколько писателей в гостевой системе ---
один читатель в эмуляторе» и позволяет отделить быструю запись сообщения от
более дорогого уведомления виртуального устройства.
